\documentclass[conference]{IEEEtran}
\IEEEoverridecommandlockouts
% The preceding line is only needed to identify funding in the first footnote. If that is unneeded, please comment it out.
\usepackage{cite}
\usepackage{amsmath,amssymb,amsfonts}
\usepackage{algorithmic}
\usepackage{graphicx}
\usepackage{textcomp}
\usepackage{xcolor}
\usepackage{verbatim}
\def\BibTeX{{\rm B\kern-.05em{\sc i\kern-.025em b}\kern-.08em
    T\kern-.1667em\lower.7ex\hbox{E}\kern-.125emX}}
\begin{document}

\title{v1.0 all experiments done. Paper waits to be completed.\\
\thanks{Identify applicable funding agency here. If none, delete this.}
}

\author{\IEEEauthorblockN{1\textsuperscript{st} Given Name Surname}
\IEEEauthorblockA{\textit{dept. name of organization (of Aff.)} \\
\textit{name of organization (of Aff.)}\\
City, Country \\
email address or ORCID}
\and
\IEEEauthorblockN{2\textsuperscript{th} Given Name Surname}
\IEEEauthorblockA{\textit{dept. name of organization (of Aff.)} \\
\textit{name of organization (of Aff.)}\\
City, Country \\
email address or ORCID}
}

\maketitle

\begin{abstract}
To be Done. 
\end{abstract}

\begin{IEEEkeywords}
DAO Governance, transaction cost economies, technology standardization, open-source governance, consortinium governance, blockchain, AI propocal
\end{IEEEkeywords}

\section{(Only-in-Draft) README}
This is a paragraph in draft to remind myself what's incomplete and what should be discussed. 

\textbf{What I have done.} I have finished most of the experimental work, including data collection, processing, analysis, metrics computation, visualization. The \textit{Research Design and Methods} and \textit{Results} parts shows where was the data from and how it was processed. All outputs can be seen on GitHub repository. What stands out: 

\begin{verbatim}
output/network_compare.html
output/timeline_erc8004.html
output/bipartite_a2a.html
output/bipartite_erc8004.html
\end{verbatim}

\textbf{What I have NOT done.} The theoretical framework remains to be completed. This is my bad: at this point, I'm not informed about the whole framework. I don't have at least 95\% confidencd on my implications, from which theory I should start, and where extend. Some more time should be finished to complete this time. Therefore, the \textit{Introduction, Theoretical Background} and \textit{Discussion} parts are only preliminary drafts.

Below are thing I should work on. 

References (to be verified, enriched and improved).
\begin{itemize}
    \item Transaction Cost Economies. (Cite classic papers? I don't think for this theories I need very new papers)
    \item Why people work without contracts? (E.g. commons-based peer production. Reputation externality. )
\end{itemize}

Some details (e.g. authors, cohen's kappa). Language polish and enrich.

My concerns (to be discussed):
\begin{itemize}
    \item Is it too old to start with and extend from the transaction cost theory? (Though it's from a Nobel prize winner.) Several more theories have been developed and they are good reference materials. 
    \item Do I need to talk about limitations of this paper?
    \item Should the paper be consistent with the RQ1 I put forward in the extended abstract?
    \item How specific should I be on introducing DAO?
\end{itemize}

\section{Introduction}
This paper compares two cases: ERC-8004 and Google A2A. Both are AI protocals, one extends from the other, but different governance model. 

Theory we start from. What we curious about. Why we compare these two cases? Consortium Governance VS DAO Governance. DAO governance: beyond markets, hybrid and hierarchy. 

The main research question.

The results and discussions.

\section{Theoretical Background}
Where we start and where we want to arrive. 
\subsection{Transaction Cost and Governance Modes}
Coase's 'nature of the firm'.
\subsection{Open-Source Governance}
How open-source governnace works and several open-source theories.
\subsection{DAO Governance as a Novel Governance Mode}
DAO and how it works. Why we choose Ethereum and ERC-8004. 

\section{Research Design and Methods}
Design and methodology.
\subsection{Comparative Case Study Design}
Why case study matters, and why choosing the two cases make sense. Some methods to follow about comparative case studies.

\subsection{Data Collection}
SHA-256 was used for data provance. 

Data are collected from three sources for each case. For ERC-8004: (1) 113 posts from the Ethereum Magicians forum thread (topic 25098), scraped via the Discourse JSON API; and (2) 36 GitHub records from 9 core lifecycle pull requests (\#1170, \#1244, \#1248, \#1458, \#1462, \#1470, \#1472, \#1477, \#1488) that directly modified \texttt{ERCS/erc-8004.md}.

For Google A2A: (1) 3,104 issue and issue comment records; (2) 1,955 pull request and review comment records; and (3) 822 GitHub Discussion records scraped via the GitHub GraphQL API.

Total records before filtering: 6,030. After removing records with fewer than 20 characters of text (bot messages, CI notifications) and bot authors: 4,323 (ERC-8004: 142; Google A2A: 4,181).


\subsection{LLM-Assisted Annotation}
Minimax-M2.5 was used for annotation. Each record received four categorical labels:
\begin{itemize}
  \item \textbf{Stakeholder institution}: Google, MetaMask, Ethereum
        Foundation, Coinbase, Independent, Unknown
  \item \textbf{Argument type}: Technical, Economic, Governance-Principle, Process, Off-topic
  \item \textbf{Stance}: Support, Oppose, Modify, Neutral, Off-topic
  \item \textbf{Consensus signal}: Adopted, Rejected, Pending, N/A
\end{itemize}
Institution labels follow a three-tier provenance cascade: (1) EIP header email (highest confidence, 2 authors); (2) R07 manual investigation---a systematic review of GitHub profiles, LinkedIn, and EIP metadata for the top 109 contributors across both cases, yielding 40 institution upgrades; and (3) LLM inference (517 remaining authors). The original LLM-inferred label is preserved in a separate \texttt{institution\_lm} field for all records.

Manual verification was also conducted on all 71 people related to erc-8004 and several main contributors of Google A2A. 

\subsection{Validity}
To be completed. Not very sure if this part is necessary and how should I do it. 

\section{Results}
See $output/findings_summary.md$.
\subsection{Key Metrics}

\subsection{Argument Type Distribution}
Here put the bipartite diagrams.

\subsection{Stance Distribution}
Here put the timeline diagram.

\subsection{Institutions}
Here put the comparitive diagram.
\subsection{Votes and Consensus}

\section{Discussion}

Intepretations of the data. 

\subsection{Consensus Instead of Contracts}
The development of ERC-8004 does not involve any voting prosess, despite that it's a Ethereum protocal. In contrast, Google A2A has several votes on GitHub, which ... Several processes that requires voting of DAOs... However, ERC-8004 relies on consensus. This aligns with the \cite{Benkler}: commons-based peer production... 

\subsection{Permissionless Entry}
Any one can deploy the codes... smart contract... immutable... 
Open-source community, trustless, permissionless. Several permissionsless theories cite here. 

\subsection{Psudonymous Reputation Externalization}
This is the blockchain and Web 3 community... Several theories cite here. 

\section{Conclusion}
Motivation, methoodlogy, findings, discussions, contributions. 

\section*{References}
\begin{thebibliography}{00}
\bibitem{firm-nature} R. H. Coase, "The nature of the firm," Economica, vol. 4, no. 16, pp. 386-405, Nov. 1937.
\bibitem{Benkler} Y. Benkler, "Coase's penguin, or, Linux and 'The Nature of the Firm'," Yale Law Journal, vol. 112, no. 3, pp.369-446, Dec. 2002.
\bibitem{consensus-governance} T. S. Simcoe, "Standard setting committees: Consensus governance for shared technology platforms," Management Science, vol. 58, no. 4, pp. 702-720, Apr. 2012.
\end{thebibliography}
\vspace{12pt}

\end{document}
